\documentclass{beamer}\usepackage[]{graphicx}\usepackage[]{color}
%% maxwidth is the original width if it is less than linewidth
%% otherwise use linewidth (to make sure the graphics do not exceed the margin)
\makeatletter
\def\maxwidth{ %
  \ifdim\Gin@nat@width>\linewidth
    \linewidth
  \else
    \Gin@nat@width
  \fi
}
\makeatother

\definecolor{fgcolor}{rgb}{0.345, 0.345, 0.345}
\newcommand{\hlnum}[1]{\textcolor[rgb]{0.686,0.059,0.569}{#1}}%
\newcommand{\hlstr}[1]{\textcolor[rgb]{0.192,0.494,0.8}{#1}}%
\newcommand{\hlcom}[1]{\textcolor[rgb]{0.678,0.584,0.686}{\textit{#1}}}%
\newcommand{\hlopt}[1]{\textcolor[rgb]{0,0,0}{#1}}%
\newcommand{\hlstd}[1]{\textcolor[rgb]{0.345,0.345,0.345}{#1}}%
\newcommand{\hlkwa}[1]{\textcolor[rgb]{0.161,0.373,0.58}{\textbf{#1}}}%
\newcommand{\hlkwb}[1]{\textcolor[rgb]{0.69,0.353,0.396}{#1}}%
\newcommand{\hlkwc}[1]{\textcolor[rgb]{0.333,0.667,0.333}{#1}}%
\newcommand{\hlkwd}[1]{\textcolor[rgb]{0.737,0.353,0.396}{\textbf{#1}}}%

\usepackage{framed}
\makeatletter
\newenvironment{kframe}{%
 \def\at@end@of@kframe{}%
 \ifinner\ifhmode%
  \def\at@end@of@kframe{\end{minipage}}%
  \begin{minipage}{\columnwidth}%
 \fi\fi%
 \def\FrameCommand##1{\hskip\@totalleftmargin \hskip-\fboxsep
 \colorbox{shadecolor}{##1}\hskip-\fboxsep
     % There is no \\@totalrightmargin, so:
     \hskip-\linewidth \hskip-\@totalleftmargin \hskip\columnwidth}%
 \MakeFramed {\advance\hsize-\width
   \@totalleftmargin\z@ \linewidth\hsize
   \@setminipage}}%
 {\par\unskip\endMakeFramed%
 \at@end@of@kframe}
\makeatother

\definecolor{shadecolor}{rgb}{.97, .97, .97}
\definecolor{messagecolor}{rgb}{0, 0, 0}
\definecolor{warningcolor}{rgb}{1, 0, 1}
\definecolor{errorcolor}{rgb}{1, 0, 0}
\newenvironment{knitrout}{}{} % an empty environment to be redefined in TeX

\usepackage{alltt} 
% \usepackage{graphicx}
\usepackage{graphics}
\usepackage[T1]{fontenc}
\usepackage{verbatim}
\usepackage{etoolbox}
\makeatletter
\preto{\@verbatim}{\topsep=-6pt \partopsep=-6pt }
\makeatother
%\usepackage{fix-cm}
\setbeamercovered{transparent}


\renewcommand{\ni}{\noindent}


% \SweaveOpts{cache=TRUE, background="white"}


\title[ 1-Intro]{1 - Intro to R}
\subtitle{Basics}

\author[ K. Maurer, S. VanderPlas, E. Hare]{Karsten Maurer, Eric Hare, and Susan VanderPlas}
\date{February 12, 2014}
\institute[ISU]{Iowa State University}
\IfFileExists{upquote.sty}{\usepackage{upquote}}{}

\begin{document}

\begin{frame}
    \maketitle
\end{frame}





\begin{frame}[fragile]
    \frametitle{Overgrown Calculator}
    \framesubtitle{Basic mathematical operators}
    
\begin{knitrout}\footnotesize
\definecolor{shadecolor}{rgb}{1, 1, 1}\color{fgcolor}\begin{kframe}
\begin{alltt}
\hlcom{# Addition}
\hlnum{2} \hlopt{+} \hlnum{2}
\end{alltt}
\begin{verbatim}
## [1] 4
\end{verbatim}
\begin{alltt}
\hlcom{# Subtraction}
\hlnum{15} \hlopt{-} \hlnum{7}
\end{alltt}
\begin{verbatim}
## [1] 8
\end{verbatim}
\begin{alltt}
\hlcom{# Multiplication}
\hlnum{109} \hlopt{*} \hlnum{23452}
\end{alltt}
\begin{verbatim}
## [1] 2556268
\end{verbatim}
\begin{alltt}
\hlcom{# Division}
\hlnum{3}\hlopt{/}\hlnum{7}
\end{alltt}
\begin{verbatim}
## [1] 0.4286
\end{verbatim}
\end{kframe}
\end{knitrout}

\end{frame}

\begin{frame}[fragile]
    \frametitle{Overgrown Calculator}
    \framesubtitle{Basic mathematical operators}
\begin{knitrout}\footnotesize
\definecolor{shadecolor}{rgb}{1, 1, 1}\color{fgcolor}\begin{kframe}
\begin{alltt}
\hlcom{# Integer division}
\hlnum{7}\hlopt{%/%}\hlnum{2}
\end{alltt}
\begin{verbatim}
## [1] 3
\end{verbatim}
\begin{alltt}
\hlcom{# Modulo operator (Remainder)}
\hlnum{7}\hlopt{%%}\hlnum{2}
\end{alltt}
\begin{verbatim}
## [1] 1
\end{verbatim}
\begin{alltt}
\hlcom{# Powers}
\hlnum{1.5}\hlopt{^}\hlnum{3}
\end{alltt}
\begin{verbatim}
## [1] 3.375
\end{verbatim}
\end{kframe}
\end{knitrout}

\end{frame}

\begin{frame}[fragile]
    \frametitle{Overgrown Calculator}
    \framesubtitle{Other functions}
    \begin{itemize}
        \item Exponentiation
        \begin{itemize}\item exp(x) \end{itemize}
        \item Logarithms
        \begin{itemize}\item log(x, base = 2) \end{itemize}
        \item Trigonometric functions
            \begin{itemize}
                \item sin(x)
                \item asin(x)
                \item cos(x)
                \item tan(x)
            \end{itemize}
        \item Variables
            \begin{itemize}
                \item MyAge <- 25
                \item Height = "5ft4in"
            \end{itemize}
    \end{itemize}
\end{frame}

% Go to code where we play around with tips data

% operating on faithful.  Find min

\begin{frame}[fragile]
    \frametitle{Variables}
    \framesubtitle{Variable Names}
    \begin{itemize}
        \item Variable names can't start with a number
        \item R is case-sensitive
        \item Some common letters are used internally by R and should be avoided as variable names (c, q, t, C, D, F, T, I)
        \item There are reserved words that R won't let you use for variable names. (for, in, while, if, else, repeat, break, next)
        \item R will let you use the name of a predefined function.  Try not to overwrite those though!
        \item ?make.names
    \end{itemize}
\end{frame}

%\begin{frame}
%    \frametitle{Basics}
%    \begin{itemize}
%        \item Basic algebra is the same
%        \item Use 2*x not 2x
%        \item Use 2\textasciicircum p instead of $2^p$
%        \item When storing into a variable use "<-" instead of "=".
%        \item Everything in R is a vector
%        \item Index a vector using []
%    \end{itemize}
%\end{frame}

\begin{frame}[fragile]
    \frametitle{Basics}
    \framesubtitle{Examples}

\begin{knitrout}\footnotesize
\definecolor{shadecolor}{rgb}{1, 1, 1}\color{fgcolor}\begin{kframe}
\begin{alltt}
\hlkwd{data}\hlstd{(tips,} \hlkwc{package} \hlstd{=} \hlstr{"reshape2"}\hlstd{)}
\hlstd{(bill} \hlkwb{<-} \hlkwd{head}\hlstd{(tips}\hlopt{$}\hlstd{total_bill))}
\end{alltt}
\begin{verbatim}
## [1] 16.99 10.34 21.01 23.68 24.59 25.29
\end{verbatim}
\begin{alltt}
\hlkwd{log}\hlstd{(bill)}
\end{alltt}
\begin{verbatim}
## [1] 2.833 2.336 3.045 3.165 3.202 3.230
\end{verbatim}
\begin{alltt}
\hlstd{bill[}\hlnum{2}\hlstd{]}
\end{alltt}
\begin{verbatim}
## [1] 10.34
\end{verbatim}
\begin{alltt}
\hlkwd{sum}\hlstd{(bill)}
\end{alltt}
\begin{verbatim}
## [1] 121.9
\end{verbatim}
\begin{alltt}
\hlstd{(bill_in_euros} \hlkwb{<-} \hlstd{bill} \hlopt{*} \hlnum{0.7982}\hlstd{)}
\end{alltt}
\begin{verbatim}
## [1] 13.561  8.253 16.770 18.901 19.628 20.186
\end{verbatim}
\end{kframe}
\end{knitrout}


\end{frame}

\begin{frame}
    \frametitle{Getting Help}
    \begin{itemize}
        \item help.start()
        \item help(command)
        \item ?command
        \item help.search("command")
%        \item apropos("pattern")
        \item View the code! Type function name or use getAnywhere("fun")
        \item Google ("R + statistics + ...")
%        \item Rseek
        \item StackOverflow
    \end{itemize}
Getting Out!
    \begin{itemize}
        \item q()
    \end{itemize}
    
\end{frame}

\begin{frame}
    \frametitle{R Reference Card}
Download the reference card from \url{http://cran.r-project.org/doc/contrib/Short-refcard.pdf}

Having this open or printed off and near you while working is helpful.
\end{frame}

\begin{frame}[fragile]
    \frametitle{Your Turn}
    \begin{itemize}

        \item Find out how many rows and columns the `iris' data set has.  Figure out at least 2 ways to do this.\\ 
        Hint: "Variable Information" section on the first page of the reference card!
        \item Use \texttt{rep} to construct the following vector:  1 1 2 2 3 3 4 4 5 5 \\
        Hint: "Data Creation" section of the reference card
        \item Use \texttt{rep} to construct this vector: 1 2 3 4 5 1 2 3 4 5 1 2 3 4 5
    \end{itemize}
\end{frame}

\begin{frame}[fragile]
\frametitle{Vectors}
\begin{itemize}
  \item A vector is a list of values that are all the same type:
\begin{knitrout}\footnotesize
\definecolor{shadecolor}{rgb}{1, 1, 1}\color{fgcolor}\begin{kframe}
\begin{alltt}
\hlstd{a} \hlkwb{<-} \hlnum{1}\hlopt{:}\hlnum{6}
\hlstd{a}
\end{alltt}
\begin{verbatim}
## [1] 1 2 3 4 5 6
\end{verbatim}
\end{kframe}
\end{knitrout}

  \item We can refer to an element of a vector by using its index: first, second, third:
\begin{knitrout}\footnotesize
\definecolor{shadecolor}{rgb}{1, 1, 1}\color{fgcolor}\begin{kframe}
\begin{alltt}
\hlstd{b} \hlkwb{<-} \hlkwd{c}\hlstd{(}\hlnum{3}\hlstd{,} \hlnum{4}\hlstd{,} \hlnum{5}\hlstd{,} \hlnum{1}\hlstd{,} \hlnum{6}\hlstd{,} \hlnum{2}\hlstd{)}
\hlstd{b[}\hlnum{3}\hlstd{]}
\end{alltt}
\begin{verbatim}
## [1] 5
\end{verbatim}
\begin{alltt}
\hlstd{b[}\hlnum{6}\hlstd{]}
\end{alltt}
\begin{verbatim}
## [1] 2
\end{verbatim}
\end{kframe}
\end{knitrout}

\end{itemize}
\end{frame}

\begin{frame}[fragile]
    \frametitle{Data Frames}
    \framesubtitle{Quick Intro}
    \begin{itemize}
        \item `tips' is a data frame.
        \item Data frames hold data sets
        \item Not every column need be the same type - like an Excel spreadsheet
        \item Each column in a data frame is a vector - so each column needs to have values that are all the same type.
        \item We can access different columns using the \$ operator.
%         \item We index data frames by using square brackets: \\
%         \verb|dataFrameName[row index, column index]|
    \end{itemize}
\end{frame}


\begin{frame}[fragile]
    \frametitle{Vectors}
    \begin{itemize}
        \item As mentioned earlier, almost everying in R is a vector.
        %\item There are multiple ways to make a vector from scratch
    \end{itemize}
\begin{knitrout}\footnotesize
\definecolor{shadecolor}{rgb}{1, 1, 1}\color{fgcolor}\begin{kframe}
\begin{alltt}
\hlstd{tip} \hlkwb{<-} \hlstd{tips}\hlopt{$}\hlstd{tip}
\hlstd{bill} \hlkwb{<-} \hlstd{tips}\hlopt{$}\hlstd{total_bill}
\hlkwd{head}\hlstd{(tip)}
\end{alltt}
\begin{verbatim}
## [1] 1.01 1.66 3.50 3.31 3.61 4.71
\end{verbatim}
\end{kframe}
\end{knitrout}

\end{frame}

%% Rework with tips

\begin{frame}[fragile]
    \frametitle{Indexing}
Sometimes we want to just grab part of a vector/matrix/dataframe.  

Vectors in R are 1-indexed - that is, we count 1, 2, 3.
\begin{knitrout}\footnotesize
\definecolor{shadecolor}{rgb}{1, 1, 1}\color{fgcolor}\begin{kframe}
\begin{alltt}
\hlkwd{head}\hlstd{(tip)}
\end{alltt}
\begin{verbatim}
## [1] 1.01 1.66 3.50 3.31 3.61 4.71
\end{verbatim}
\begin{alltt}
\hlstd{tip[}\hlnum{1}\hlstd{]}
\end{alltt}
\begin{verbatim}
## [1] 1.01
\end{verbatim}
\begin{alltt}
\hlkwd{c}\hlstd{(}\hlnum{1}\hlstd{,} \hlnum{3}\hlstd{,} \hlnum{5}\hlstd{)}
\end{alltt}
\begin{verbatim}
## [1] 1 3 5
\end{verbatim}
\begin{alltt}
\hlstd{tip[}\hlkwd{c}\hlstd{(}\hlnum{1}\hlstd{,} \hlnum{3}\hlstd{,} \hlnum{5}\hlstd{)]}
\end{alltt}
\begin{verbatim}
## [1] 1.01 3.50 3.61
\end{verbatim}
\begin{alltt}
\hlstd{tip[}\hlnum{1}\hlopt{:}\hlnum{5}\hlstd{]}
\end{alltt}
\begin{verbatim}
## [1] 1.01 1.66 3.50 3.31 3.61
\end{verbatim}
\end{kframe}
\end{knitrout}

\end{frame}

\begin{frame}[fragile]
    \frametitle{Helpful functions}
\begin{knitrout}\footnotesize
\definecolor{shadecolor}{rgb}{1, 1, 1}\color{fgcolor}\begin{kframe}
\begin{alltt}
\hlkwd{min}\hlstd{(tip)}
\end{alltt}
\begin{verbatim}
## [1] 1
\end{verbatim}
\begin{alltt}
\hlkwd{which.min}\hlstd{(tip)}
\end{alltt}
\begin{verbatim}
## [1] 68
\end{verbatim}
\begin{alltt}
\hlstd{tip[}\hlkwd{which.min}\hlstd{(tip)]}
\end{alltt}
\begin{verbatim}
## [1] 1
\end{verbatim}
\begin{alltt}
\hlstd{size} \hlkwb{<-} \hlstd{tips}\hlopt{$}\hlstd{size}
\hlkwd{which.min}\hlstd{(size)}
\end{alltt}
\begin{verbatim}
## [1] 68
\end{verbatim}
\begin{alltt}
\hlkwd{which}\hlstd{(size} \hlopt{==} \hlkwd{min}\hlstd{(size))}
\end{alltt}
\begin{verbatim}
## [1]  68  83 112 223
\end{verbatim}
\end{kframe}
\end{knitrout}


\end{frame}

\begin{frame}[fragile]
    \frametitle{Vectors}
        Common operations and functions work on every element of a vector and return a new vector
\begin{knitrout}\footnotesize
\definecolor{shadecolor}{rgb}{1, 1, 1}\color{fgcolor}\begin{kframe}
\begin{alltt}
\hlkwd{head}\hlstd{(tip} \hlopt{*} \hlnum{100}\hlstd{)}
\end{alltt}
\begin{verbatim}
## [1] 101 166 350 331 361 471
\end{verbatim}
\begin{alltt}
\hlkwd{head}\hlstd{(bill} \hlopt{*} \hlnum{0.18}\hlstd{)}
\end{alltt}
\begin{verbatim}
## [1] 3.058 1.861 3.782 4.262 4.426 4.552
\end{verbatim}
\begin{alltt}
\hlkwd{head}\hlstd{(}\hlkwd{round}\hlstd{(bill} \hlopt{*} \hlnum{0.18}\hlstd{,} \hlnum{2}\hlstd{))}
\end{alltt}
\begin{verbatim}
## [1] 3.06 1.86 3.78 4.26 4.43 4.55
\end{verbatim}
\begin{alltt}
\hlstd{rate} \hlkwb{<-} \hlstd{tip}\hlopt{/}\hlstd{bill}
\hlkwd{head}\hlstd{(rate)}
\end{alltt}
\begin{verbatim}
## [1] 0.05945 0.16054 0.16659 0.13978 0.14681 0.18624
\end{verbatim}
\end{kframe}
\end{knitrout}

\end{frame}


\begin{frame}
    \frametitle{Logical Values}
    \begin{itemize}
        \item R has built in support for logical values
        \item TRUE and FALSE are built in.  T (for TRUE) and F (for FALSE) are supported but can be modified
        \item Logicals can result from a comparison using
        \begin{itemize}
            \item $<$
            \item $>$
            \item $<=$
            \item $>=$
            \item $==$
            \item $!=$
        \end{itemize}
    \end{itemize}
\end{frame}

\begin{frame}[fragile]
    \frametitle{Indexing}
We can index using logical values as well
\begin{knitrout}\footnotesize
\definecolor{shadecolor}{rgb}{1, 1, 1}\color{fgcolor}\begin{kframe}
\begin{alltt}
\hlstd{x} \hlkwb{<-} \hlkwd{c}\hlstd{(}\hlnum{2}\hlstd{,} \hlnum{3}\hlstd{,} \hlnum{5}\hlstd{,} \hlnum{7}\hlstd{)}
\hlstd{x[}\hlkwd{c}\hlstd{(T, F, F, T)]}
\end{alltt}
\begin{verbatim}
## [1] 2 7
\end{verbatim}
\begin{alltt}
\hlstd{x} \hlopt{>} \hlnum{3.5}
\end{alltt}
\begin{verbatim}
## [1] FALSE FALSE  TRUE  TRUE
\end{verbatim}
\begin{alltt}
\hlstd{x[x} \hlopt{>} \hlnum{3.5}\hlstd{]}
\end{alltt}
\begin{verbatim}
## [1] 5 7
\end{verbatim}
\begin{alltt}
\hlstd{x[}\hlkwd{c}\hlstd{(T, F)]}
\end{alltt}
\begin{verbatim}
## [1] 2 5
\end{verbatim}
\end{kframe}
\end{knitrout}

\end{frame}

\begin{frame}[fragile]
    \frametitle{Logical Values}
    \framesubtitle{Examples}
\begin{knitrout}\footnotesize
\definecolor{shadecolor}{rgb}{1, 1, 1}\color{fgcolor}\begin{kframe}
\begin{alltt}
\hlkwd{head}\hlstd{(rate)}
\end{alltt}
\begin{verbatim}
## [1] 0.05945 0.16054 0.16659 0.13978 0.14681 0.18624
\end{verbatim}
\begin{alltt}
\hlstd{sad_tip} \hlkwb{<-} \hlstd{rate} \hlopt{<} \hlnum{0.1}
\hlkwd{head}\hlstd{(sad_tip)}
\end{alltt}
\begin{verbatim}
## [1]  TRUE FALSE FALSE FALSE FALSE FALSE
\end{verbatim}
\begin{alltt}
\hlstd{rate[sad_tip]}
\end{alltt}
\begin{verbatim}
##  [1] 0.05945 0.07180 0.07893 0.05680 0.09936 0.05643 0.09553
##  [8] 0.07862 0.07296 0.08147 0.09984 0.09453 0.07718 0.07398
## [15] 0.06566 0.09560 0.09001 0.07746 0.08364 0.06653 0.08527
## [22] 0.08330 0.07937 0.03564 0.07358 0.08822 0.09820
\end{verbatim}
\end{kframe}
\end{knitrout}

\end{frame}


\begin{frame}[fragile]
    \frametitle{Indexing}
We can modify subsets of vectors
\begin{knitrout}\footnotesize
\definecolor{shadecolor}{rgb}{1, 1, 1}\color{fgcolor}\begin{kframe}
\begin{alltt}
\hlstd{x} \hlkwb{<-} \hlstd{bill[}\hlnum{1}\hlopt{:}\hlnum{5}\hlstd{]}
\hlstd{x}
\end{alltt}
\begin{verbatim}
## [1] 16.99 10.34 21.01 23.68 24.59
\end{verbatim}
\begin{alltt}
\hlstd{x[}\hlnum{1}\hlstd{]}
\end{alltt}
\begin{verbatim}
## [1] 16.99
\end{verbatim}
\begin{alltt}
\hlstd{x[}\hlnum{1}\hlstd{]} \hlkwb{<-} \hlnum{20}
\hlstd{x}
\end{alltt}
\begin{verbatim}
## [1] 20.00 10.34 21.01 23.68 24.59
\end{verbatim}
\end{kframe}
\end{knitrout}

\end{frame}

\small
\begin{frame}[fragile]
    \frametitle{Vectors}
Elements of a vector all must be the same type.

\begin{knitrout}\footnotesize
\definecolor{shadecolor}{rgb}{1, 1, 1}\color{fgcolor}\begin{kframe}
\begin{alltt}
\hlkwd{head}\hlstd{(rate)}
\end{alltt}
\begin{verbatim}
## [1] 0.05945 0.16054 0.16659 0.13978 0.14681 0.18624
\end{verbatim}
\begin{alltt}
\hlstd{rate[sad_tip]} \hlkwb{<-} \hlstr{":-("}
\hlkwd{head}\hlstd{(rate)}
\end{alltt}
\begin{verbatim}
## [1] ":-("               "0.160541586073501"
## [3] "0.166587339362208" "0.139780405405405"
## [5] "0.146807645384303" "0.186239620403321"
\end{verbatim}
\end{kframe}
\end{knitrout}


All of the items in \texttt{rate} are now strings! That is sad!
\end{frame}
\normalsize


\begin{frame}[fragile]
    \frametitle{Data types}
    \begin{itemize}
        \item Can use `mode' or `class' to find out information about variables
        \item `str` is useful to find information about the structure of your data
        \item Many data types
        \begin{itemize}
            \item integer
            \item numeric
            \item character
            \item Date
            \item factor
        \end{itemize}
    \end{itemize}


\begin{knitrout}\footnotesize
\definecolor{shadecolor}{rgb}{1, 1, 1}\color{fgcolor}\begin{kframe}
\begin{alltt}
\hlkwd{str}\hlstd{(tips)}
\end{alltt}
\begin{verbatim}
## 'data.frame':	244 obs. of  7 variables:
##  $ total_bill: num  17 10.3 21 23.7 24.6 ...
##  $ tip       : num  1.01 1.66 3.5 3.31 3.61 4.71 2 3.12 1.96 3.23 ...
##  $ sex       : Factor w/ 2 levels "Female","Male": 1 2 2 2 1 2 2 2 2 2 ...
##  $ smoker    : Factor w/ 2 levels "No","Yes": 1 1 1 1 1 1 1 1 1 1 ...
##  $ day       : Factor w/ 4 levels "Fri","Sat","Sun",..: 3 3 3 3 3 3 3 3 3 3 ...
##  $ time      : Factor w/ 2 levels "Dinner","Lunch": 1 1 1 1 1 1 1 1 1 1 ...
##  $ size      : int  2 3 3 2 4 4 2 4 2 2 ...
\end{verbatim}
\end{kframe}
\end{knitrout}

\end{frame}

\begin{frame}[fragile]
    \frametitle{Data types}
    \framesubtitle{Converting between types}
Can convert between types using \verb|as.____|

\begin{knitrout}\footnotesize
\definecolor{shadecolor}{rgb}{1, 1, 1}\color{fgcolor}\begin{kframe}
\begin{alltt}
\hlstd{size} \hlkwb{<-} \hlkwd{head}\hlstd{(tips}\hlopt{$}\hlstd{size)}
\hlstd{size}
\end{alltt}
\begin{verbatim}
## [1] 2 3 3 2 4 4
\end{verbatim}
\begin{alltt}
\hlkwd{as.character}\hlstd{(size)}
\end{alltt}
\begin{verbatim}
## [1] "2" "3" "3" "2" "4" "4"
\end{verbatim}
\begin{alltt}
\hlkwd{as.numeric}\hlstd{(}\hlstr{"2"}\hlstd{)}
\end{alltt}
\begin{verbatim}
## [1] 2
\end{verbatim}
\begin{alltt}
\hlkwd{as.factor}\hlstd{(size)}
\end{alltt}
\begin{verbatim}
## [1] 2 3 3 2 4 4
## Levels: 2 3 4
\end{verbatim}
\end{kframe}
\end{knitrout}

\end{frame}

\begin{frame}[fragile]
    \frametitle{Data types}
    \framesubtitle{Summary}

\begin{knitrout}\footnotesize
\definecolor{shadecolor}{rgb}{1, 1, 1}\color{fgcolor}\begin{kframe}
\begin{alltt}
\hlstd{size} \hlkwb{<-} \hlstd{tips}\hlopt{$}\hlstd{size}
\hlkwd{summary}\hlstd{(size)}
\end{alltt}
\begin{verbatim}
##    Min. 1st Qu.  Median    Mean 3rd Qu.    Max. 
##    1.00    2.00    2.00    2.57    3.00    6.00
\end{verbatim}
\begin{alltt}
\hlkwd{summary}\hlstd{(}\hlkwd{as.character}\hlstd{(size))}
\end{alltt}
\begin{verbatim}
##    Length     Class      Mode 
##       244 character character
\end{verbatim}
\begin{alltt}
\hlkwd{summary}\hlstd{(}\hlkwd{as.factor}\hlstd{(size))}
\end{alltt}
\begin{verbatim}
##   1   2   3   4   5   6 
##   4 156  38  37   5   4
\end{verbatim}
\end{kframe}
\end{knitrout}

\end{frame}

\begin{frame}
    \frametitle{Functions}
    \begin{itemize}
        \item Typical format: foo(x, n = length(x), ...)
        \item Some parameters have defaults set for you
        \item ... is special.  This passes along extra parameters to functions used inside the function
        \item Use args(foo) to view arguments for functions.
    \end{itemize}
\end{frame}


%\begin{frame}[fragile]
%    \frametitle{Vectors}
%% latex table generated in R 2.15.0 by xtable 1.7-0 package
%% Fri May 25 12:27:08 2012
%\begin{table}[ht]
%\begin{center}
%\begin{tabular}{|r|rrr|}
%  \hline
%ID & field & plot & amount \\ 
%  \hline
%1 &   1 &   1 & 10.0 \\ 
%  2 &   1 &   2 & 15.0 \\ 
%  3 &   1 &   3 & 20.0 \\ 
%  4 &   1 &   4 & 25.0 \\ 
%  5 &   2 &   1 & 10.0 \\ 
%  6 &   2 &   2 & 15.0 \\ 
%  7 &   2 &   3 & 20.0 \\ 
%  8 &   2 &   4 & 25.0 \\ 
%   \hline
%\end{tabular}
%\end{center}
%\end{table}
%\end{frame}


\begin{frame}[fragile]
    \frametitle{Basic Statistical Functions}
Using the basic functions we've learned it wouldn't be hard to compute some basic statistics.
\begin{knitrout}\footnotesize
\definecolor{shadecolor}{rgb}{1, 1, 1}\color{fgcolor}\begin{kframe}
\begin{alltt}
\hlstd{(n} \hlkwb{<-} \hlkwd{length}\hlstd{(tip))}
\end{alltt}
\begin{verbatim}
## [1] 244
\end{verbatim}
\begin{alltt}
\hlstd{(meantip} \hlkwb{<-} \hlkwd{sum}\hlstd{(tip)}\hlopt{/}\hlstd{n)}
\end{alltt}
\begin{verbatim}
## [1] 2.998
\end{verbatim}
\begin{alltt}
\hlstd{(standdev} \hlkwb{<-} \hlkwd{sqrt}\hlstd{(}\hlkwd{sum}\hlstd{((tip} \hlopt{-} \hlstd{meantip)}\hlopt{^}\hlnum{2}\hlstd{)}\hlopt{/}\hlstd{(n} \hlopt{-} \hlnum{1}\hlstd{)))}
\end{alltt}
\begin{verbatim}
## [1] 1.384
\end{verbatim}
\end{kframe}
\end{knitrout}

But we don't have to.
\end{frame}


\begin{frame}[fragile]
    \frametitle{Basic Statistical Functions}
\begin{knitrout}\footnotesize
\definecolor{shadecolor}{rgb}{1, 1, 1}\color{fgcolor}\begin{kframe}
\begin{alltt}
\hlkwd{mean}\hlstd{(tip)}
\end{alltt}
\begin{verbatim}
## [1] 2.998
\end{verbatim}
\begin{alltt}
\hlkwd{sd}\hlstd{(tip)}
\end{alltt}
\begin{verbatim}
## [1] 1.384
\end{verbatim}
\begin{alltt}
\hlkwd{summary}\hlstd{(tip)}
\end{alltt}
\begin{verbatim}
##    Min. 1st Qu.  Median    Mean 3rd Qu.    Max. 
##    1.00    2.00    2.90    3.00    3.56   10.00
\end{verbatim}
\begin{alltt}
\hlkwd{quantile}\hlstd{(tip,} \hlkwd{c}\hlstd{(}\hlnum{0.025}\hlstd{,} \hlnum{0.975}\hlstd{))}
\end{alltt}
\begin{verbatim}
##  2.5% 97.5% 
## 1.176 6.462
\end{verbatim}
\end{kframe}
\end{knitrout}

\end{frame}

\begin{frame}
    \frametitle{Distributions}
    \begin{itemize}
        \item R has a lot of distributions built in.
        \item We can typically obtain:
        \begin{itemize}
            \item Density value
            \item CDF value
            \item Inverse CDF value (percentiles)
            \item Random deviate
        \end{itemize}
        \item Normal (\texttt{norm}), Chi-square (\texttt{chisq}), F (\texttt{f}), T (\texttt{t}), Cauchy (\texttt{cauchy}), Poisson (\texttt{pois}), Binomial (\texttt{binom}), Negative Binomial (\texttt{nbinom}), Gamma (\texttt{gamma}), ..., lots more!
        \item \texttt{library(help = stats)}
    \end{itemize}
\end{frame}

\begin{frame}[fragile]
    \frametitle{Distributions}
    \framesubtitle{Examples}
\begin{knitrout}\footnotesize
\definecolor{shadecolor}{rgb}{1, 1, 1}\color{fgcolor}\begin{kframe}
\begin{alltt}
\hlcom{# Density - prepend with 'd'}
\hlkwd{dnorm}\hlstd{(}\hlnum{0}\hlstd{,} \hlkwc{mean} \hlstd{=} \hlnum{0}\hlstd{,} \hlkwc{sd} \hlstd{=} \hlnum{1}\hlstd{)}
\end{alltt}
\begin{verbatim}
## [1] 0.3989
\end{verbatim}
\begin{alltt}
\hlcom{# CDF - prepend with 'p'}
\hlkwd{pnorm}\hlstd{(}\hlnum{2}\hlstd{)}
\end{alltt}
\begin{verbatim}
## [1] 0.9772
\end{verbatim}
\begin{alltt}
\hlcom{# Inverse CDF - prepend with 'q'}
\hlkwd{qnorm}\hlstd{(}\hlnum{0.975}\hlstd{)}
\end{alltt}
\begin{verbatim}
## [1] 1.96
\end{verbatim}
\begin{alltt}
\hlcom{# Random deviates - prepend with 'r'}
\hlkwd{rnorm}\hlstd{(}\hlnum{5}\hlstd{)}
\end{alltt}
\begin{verbatim}
## [1] -0.23205 -1.54898  1.32022 -0.42632 -0.04595
\end{verbatim}
\end{kframe}
\end{knitrout}

\end{frame}

%\begin{frame}[fragile]
%    \frametitle{Distributions}
%<<dDist, echo=FALSE, out.width=".8\\linewidth", fig.align="center">>=
%x <- seq(-3, 3, .001)
%y <- dnorm(x)
%
%plot(x, y, type = "l")
%
%id <- x <= qnorm(.95)
%polygon(c(x[id], rev(x[id])), c(y[id], rep(0, sum(id))), col = "red")
%@
%\end{frame}

%\begin{frame}
%    \frametitle{Your Turn}
%    \begin{itemize}
%        \item Find the mean of 10, 1000, and 10000 random normal observations from N(0,1)
%        \item Generate 1,000,000 random observations from N(0,1), square them, and find the .95 quantile of the observed result
%        \item What is the .95 quantile from a Chi-square distribution with 1 degree of freedom?
%        \item Store into 'x': 100 random deviates from a normal with mean 5, variance of 36
%        \item Calculate a 95\% confidence interval for $\mu$ using x as your data. $\bar{x} \pm t_{.975,df=99}\frac{s}{\sqrt{n}}$
%    \end{itemize}
%\end{frame}




\begin{frame}[fragile]
    \frametitle{Comparisons and Numeric Types}
    \framesubtitle{Be careful}



\begin{knitrout}\footnotesize
\definecolor{shadecolor}{rgb}{1, 1, 1}\color{fgcolor}\begin{kframe}
\begin{alltt}
\hlnum{2} \hlopt{==} \hlnum{2}
\end{alltt}
\begin{verbatim}
## [1] TRUE
\end{verbatim}
\begin{alltt}
\hlkwd{sqrt}\hlstd{(}\hlnum{2}\hlstd{)}\hlopt{^}\hlnum{2}
\end{alltt}
\begin{verbatim}
## [1] 2
\end{verbatim}
\begin{alltt}
\hlkwd{sqrt}\hlstd{(}\hlnum{2}\hlstd{)}\hlopt{^}\hlnum{2} \hlopt{==} \hlnum{2}
\end{alltt}
\begin{verbatim}
## [1] FALSE
\end{verbatim}
\begin{alltt}
\hlkwd{options}\hlstd{(}\hlkwc{digits} \hlstd{=} \hlnum{22}\hlstd{)}
\hlkwd{sqrt}\hlstd{(}\hlnum{2}\hlstd{)}\hlopt{^}\hlnum{2}
\end{alltt}
\begin{verbatim}
## [1] 2.000000000000000444089
\end{verbatim}
\end{kframe}
\end{knitrout}

\end{frame}

\begin{frame}[fragile]
    \frametitle{Comparisons and Numeric Types}
    \framesubtitle{Be careful}
\begin{knitrout}\footnotesize
\definecolor{shadecolor}{rgb}{1, 1, 1}\color{fgcolor}\begin{kframe}
\begin{alltt}
\hlkwd{sqrt}\hlstd{(}\hlnum{2}\hlstd{)}\hlopt{^}\hlnum{2} \hlopt{==} \hlnum{2}
\end{alltt}
\begin{verbatim}
## [1] FALSE
\end{verbatim}
\begin{alltt}
\hlkwd{all.equal}\hlstd{(}\hlkwd{sqrt}\hlstd{(}\hlnum{2}\hlstd{)}\hlopt{^}\hlnum{2}\hlstd{,} \hlnum{2}\hlstd{)}
\end{alltt}
\begin{verbatim}
## [1] TRUE
\end{verbatim}
\begin{alltt}
\hlkwd{options}\hlstd{(}\hlkwc{digits} \hlstd{=} \hlnum{22}\hlstd{)}
\hlnum{2}\hlopt{^}\hlnum{64}
\end{alltt}
\begin{verbatim}
## [1] 18446744073709551616
\end{verbatim}
\begin{alltt}
\hlnum{2}\hlopt{^}\hlnum{64} \hlopt{-} \hlnum{1}
\end{alltt}
\begin{verbatim}
## [1] 18446744073709551616
\end{verbatim}
\begin{alltt}
\hlcom{# Uh oh!  The 'gmp' or 'int64' packages provides better}
\hlcom{# integer precision if needed.}
\end{alltt}
\end{kframe}
\end{knitrout}



\end{frame}

\begin{frame}[fragile]
    \frametitle{Logical Operators}
    \begin{itemize}
        \item \& (elementwise AND)
        \item | (elementwise OR)
    \end{itemize}

\begin{knitrout}\footnotesize
\definecolor{shadecolor}{rgb}{1, 1, 1}\color{fgcolor}\begin{kframe}
\begin{alltt}
\hlkwd{c}\hlstd{(T, T, F, F)} \hlopt{&} \hlkwd{c}\hlstd{(T, F, T, F)}
\end{alltt}
\begin{verbatim}
## [1]  TRUE FALSE FALSE FALSE
\end{verbatim}
\begin{alltt}
\hlkwd{c}\hlstd{(T, T, F, F)} \hlopt{|} \hlkwd{c}\hlstd{(T, F, T, F)}
\end{alltt}
\begin{verbatim}
## [1]  TRUE  TRUE  TRUE FALSE
\end{verbatim}
\begin{alltt}
\hlcom{# Which are big bills with a poor tip rate?}
\hlstd{id} \hlkwb{<-} \hlstd{(bill} \hlopt{>} \hlnum{40} \hlopt{&} \hlstd{rate} \hlopt{<} \hlnum{0.1}\hlstd{)}
\hlstd{tips[id, ]}
\end{alltt}
\begin{verbatim}
##     total_bill tip    sex smoker day   time size
## 103      44.30 2.5 Female    Yes Sat Dinner    3
## 183      45.35 3.5   Male    Yes Sun Dinner    3
## 185      40.55 3.0   Male    Yes Sun Dinner    2
\end{verbatim}
\end{kframe}
\end{knitrout}


?"\&"

\end{frame}



%\begin{frame}[fragile]
%    \frametitle{Your Turn}
%<<dYT>>=
%pat <- seq(2, 103, by = 3)
%@
%    \begin{itemize}
%        \item Let x = elements at even indices of pat
%        \item Let y = elements in pat greater than mean(pat)
%        \item Let z = The even elements in pat
%        \item Find the indices of the elements divisible by 7 (the function `which` might help)
%    \end{itemize}
%\end{frame}

\begin{frame}[fragile]
    \frametitle{Your Turn}
    Using the `diamonds' dataset from ggplot2:
    \begin{itemize}
        \item Read up on the dataset (?diamonds)
        \item Create a variable for price/carat
        \item Compare the price of "Fair" diamonds and "Ideal" diamonds using summary.  If you want to try to do a t-test to test for a difference in the mean of the prices.
        \item Plot price by carat (use qplot - go back to the motivating example for help with the syntax)
        \item Explore any interesting relationships you find.
    \end{itemize}
\end{frame}


\end{document}
